% PDF-to-TeX transcription by the assistant; not the original downloaded file.

\begin{theorem}[37.1, Fully relative Poincar\'e duality]
Let $M$ be an $n$-manifold and $K\supseteq L$ a pair of compact subsets. Assume given an $R$-orientation along $K$, with corresponding fundamental class $[M]_K$. With $p+q=n$, the map
\[
\cap[M]_K:\check H^p(K,L;R)\longrightarrow H_q(M-L,M-K;R)
\]
is an isomorphism.
\end{theorem}
We have seen that these isomorphisms are compatible; they form the rungs of the commuting ladder
\[
\xymatrix@C=8pt{
\cdots\ar[r]&\check H^{p-1}(L)\ar[r]\ar[d]^{\cap[M]_L}&\check H^p(K,L)\ar[r]\ar[d]^{\cap[M]_K}&\check H^p(K)\ar[r]\ar[d]^{\cap[M]_K}&\check H^p(L)\ar[r]\ar[d]^{\cap[M]_L}&\cdots\\
\cdots\ar[r]&H_{q+1}(M,M-L)\ar[r]&H_q(M-L,M-K)\ar[r]&H_q(M,M-K)\ar[r]&H_q(M,M-L)\ar[r]&\cdots
}
\]
To prove this theorem, we will follow the same five-step process we used to prove the Orientation Theorem 32.1. We have already prepared the Mayer--Vietoris ladder for this purpose.
\begin{proof}[Proof of Theorem 37.1]
By the top ladder and the five-lemma, we may assume $L=\varnothing$; so we want to prove that
\[
\cap[M]_K:\check H^p(K;R)\longrightarrow H_q(M,M-K;R)
\]
is an isomorphism.

(1) $M=\R^n$, $K$ a compact convex set. We claim that
\[
\check H^*(K)\xrightarrow{\cong}H^*(K).
\]
For any $\epsilon>0$, let $U_\epsilon$ denote the $\epsilon$-neighborhood of $K$,
\[
U_\epsilon=\bigcup_{x\in K}B_\epsilon(x).
\]
For any $y\in U_\epsilon$ there is a closest point in $K$, since the distance function to $y$ is continuous and bounded below on the compact set $K$ and so achieves its infimum. If $x',x''\in K$ are the same distance from $y$, then the midpoint of the segment joining $x'$ and $x''$ is closer, but lies in $K$ since $K$ is convex. So there is a unique closest point, $f(y)$. We let the listener check that $f:U_\epsilon\to K$ is continuous. It is also clear that if $i:K\to U_\epsilon$ is the inclusion then $i\circ f$ is homotopic to the identity on $Y$, by an affine homotopy.

Now let $D^n$ be a disk centered at the origin and containing the compact set $K$, and consider the commutative diagram
\[
\xymatrix{
H^p(K)\ar[r]^{\cap[\R^n]_K}&H_q(\R^n,\R^n-K)\\
H^p(D^n)\ar[u]^{\cong}\ar[r]\ar[d]_{\cong}&H_q(\R^n,\R^n-D^n)\ar[u]^{\cong}\ar[d]_{\cong}\\
H^p(*)\ar[r]&H_q(\R^n,\R^n-*).
}
\]
The groups are zero unless $p=0$, $q=n$. By naturality of the cap product, the bottom map is given by $1\mapsto1\cap[\R^n]_*$, and this is $[\R^n]_*$ since capping with $1$ is the identity, and this fundamental class is a generator of $H_n(\R^n,\R^n-*)$.

(2) $K$ a finite union of compact convex subsets of $\R^n$. This follows by induction and the five lemma applied to the Mayer--Vietoris ladder 36.2.

(3) $K$ is any compact subset of $\R^n$. This follows as before by a limit argument, using Lemmas 32.4 and 37.2.

(4) $M$ arbitrary, $K$ is a finite union of compact Euclidean subsets of $M$. This follows from (3) and Theorem 36.2.

(5) $M$ arbitrary, $K$ an arbitrary compact subset. This follows just as in the proof of Theorem 32.1.
\end{proof}
\paragraph{Referenced source material: Lemma 32.4}
\begin{lemma}[32.4]
Let $A_1\supseteq A_2\supseteq\cdots$ be a decreasing sequence of compact subsets of a space $X$, with intersection $A$. Then
\[
\varinjlim_i H_q(X,X-A_i)\xrightarrow{\cong}H_q(X,X-A).
\]
\end{lemma}
\begin{proof}
Let $\sigma:\Delta^q\to X$ be any $q$-simplex in $X-A$. The subsets $X-A_i$ form an open cover of $\im(\sigma)$, so by compactness it lies in some single $X-A_i$. This shows that
\[
\varinjlim_i S_q(X-A_i)\xrightarrow{\cong}S_q(X-A).
\]
Thus
\[
\varinjlim_i S_q(X|A_i)\xrightarrow{\cong}S_q(X|A_i)
\]
by exactness of direct limit, and the claim then follows for the same reason.
\end{proof}
\paragraph{Referenced source material: Theorem 32.1, steps (3) and (5)}
\noindent\textit{The source calls a subset of $M$ ``Euclidean'' if it lies inside some open set homeomorphic to $\R^n$.}

(3) For each $j\geq1$, let $C_j$ be a finite subset of $A$ such that
\[
A\subseteq\bigcup_{x\in C_j}B_{1/j}(x).
\]
Since any intersection of convex sets is either empty or convex,
\[
A_k=\bigcap_{j=1}^k\ \bigcup_{x\in C_j}B_{1/j}(x)
\]
is a union of finitely many convex sets, and since $A$ is closed it is the intersection of this decreasing family. So the result follows from (1), (2), and Proposition 32.3.

(5) Cover $A$ by finitely many open subsets that embed in Euclidean opens as open disks with compact closures. Their closures then form a finite cover by closed Euclidean disks $D_i$ in Euclidean opens $U_i$. For each $i$, excise the closed subset $M-U_i$ to see that
\begin{align*}
H_q(M,M-A\cap D_i)&\cong H_q(U_i,U_i-A\cap D_i)\\
&\cong H_q(\R^n,\R^n-A\cap D_i).
\end{align*}
By (4), the theorem holds for each of these. Each intersection $(A\cap D_i)\cap(A\cap D_j)$ is again a compact Euclidean subset, so the result holds for them by excision as well. The result then follows by (1).

